\section{SF-E848: arbitrary matrix extensions over a common point complete intersection}\label{SF-E848:section}

\textbf{Checked mathematical status, 2026-11-22.} The author, root,
and independent reviewer \texttt{/root/g34\_eightfolds} have read the
complete frozen proof. All mathematical reviews passed, with no repair.
The independent full proof/source report was completed and saved at
\texttt{03:19:07 UTC}. The frozen Markdown original retains its author
status. This canonical conversion has a separate whole-text check.
No numerical certificate or strict comparison claim is added here.
The matrix dimension argument was developed jointly by \texttt{/root} and
\texttt{/root/sixfolds}; this is not a claim of independent discovery by either.
Author of this written proof: \texttt{/root/sixfolds}, GPT-6 Astra Pro; inherited
effort is not exposed. The scalar argument SF-E847 is an explicitly
declared input and is also expanded here where the matrix case changes it.
The independent review and its declared dependency scopes are recorded
above. No numerical computation or new strict movement is claimed by
the present status paragraph.

\subsection{1. Statements and actual local type}\label{SF-E848:s1}

Let \(R=\mathbf C[[x_1,x_2,x_3,x_4]]\), with maximal ideal
\(\mathfrak m\), and let \(I=(f_1,f_2,f_3,f_4)\) be a proper
\(\mathfrak m\)-primary complete-intersection ideal. Write \(A=R/I\)
and \(\ell=\dim_{\mathbf C}A>0\). No restrictions are placed on the
orders, initial forms, or nonlinear coefficients of the four generators.
For an integer \(q\ge1\), let \(b=(b_{ij})_{i<j}\) be an arbitrary
six-tuple of \(q\) by \(q\) matrices over \(A\). The Koszul resolution
identifies it with the class of an actual extension

\[
 0\longrightarrow A^q\longrightarrow M_b\longrightarrow I^q
       \longrightarrow0,
 \qquad
 b\in\operatorname{Ext}^1_R(I^q,A^q)
        =\operatorname{Mat}_q(A)^6.
 \tag{MC.1}
\]

The tuple may be zero. Its six matrices need not commute, need not be
simultaneously diagonalizable, and need not have any common scalar form.
The composite \(\phi:M_b\to I^q\to R^q\) is the canonical map to the
double dual. Indeed the grade-four vanishings for \(A\) give
\(\operatorname{Hom}_R(M_b,R)=\operatorname{Hom}_R(I^q,R)=R^q\),
with generators the actual coordinate hull maps. Thus the double dual
is \(R^q\), its image is \(I^q\), and both its kernel and cokernel are
\(A^q\).

\textbf{Local theorem SF-E848a.} Put
\(e_i=\dim_{\mathbf C}\operatorname{Ext}^i_R(M_b,M_b)\) for \(i>0\).
Then

\[
 e_2\ge(2q^2+8)\ell.
 \tag{MC.2}
\]

In particular, this bound is valid for every matrix tuple, including a
split extension. All positive Ext groups have finite length because
\(M_b\) is free off the closed point. The proof uses duality only on
a finite-supported hull-cone complex. It does not assert ordinary local
self-duality for an endomorphism complex of positive rank.

For the global statement let \(X\) be a complex abelian fourfold. Let
\(E\) be a coherent sheaf whose canonical evaluation into
\(H=E^{\vee\vee}\) has locally free hull and finite kernel and cokernel:

\[
 0\longrightarrow T\longrightarrow E\longrightarrow H
        \longrightarrow Q\longrightarrow0.
 \tag{MC.3}
\]

Suppose that at one point \(p\) its actual completed module has the form

\[
 \widehat E_p\simeq M_b\oplus R^s,
 \qquad b\in\operatorname{Mat}_m(A)^6,
 \quad m>0,\quad s\ge0,
 \tag{MC.4}
\]

with a single common complete-intersection algebra \(A\) for all
torsion and quotient copies. This is an isomorphism of actual modules,
not an equality in a Grothendieck group.

\textbf{Global conditional theorem SF-E848b.} Fix an effective coherent finite
translation/tensor action \(G\) in the supplied DEG\_L S21 framework.
Assume there is a class \(\xi\in HH^2(X)\) whose holomorphic bivector
component is nonzero and such that the actual evaluation
\(\rho_E(\xi)\) vanishes in ordinary \(\operatorname{Ext}^2_X(E,E)\).
Then the full semiregularity map on invariant Ext is not injective.
More precisely its kernel has dimension at least \(10\ell-6>0\).

\textbf{Balanced-RM consequence SF-E848c.} Suppose \(X\) has the genuine
real-multiplication structure and compatible polarization
\(\theta=\theta_++\theta_-\) of the seed problem, and

\[
 e^B\operatorname{ch}(E)
   =r\left(1-\nu\theta_+^2/2\right)
        \left(1-\nu\theta_-^2/2\right),
 \qquad r>0,\quad\nu>0,
 \tag{MC.5}
\]

where \(B\) is a rational algebraic divisor. A sheaf satisfying
(MC.3)--(MC.4) cannot be weakly \(G\)-semiregular for any such specified
action. The proof will obtain the vanishing required in SF-E848b from
the balanced character and the hypothetical injectivity of the
semiregularity map. It does not require a balanced character on any
individual stabilizer summand.

For the identity action alone, the local bound (MC.2) already gives
failure of semiregularity for (MC.4), without (MC.5) or the hypothesis
on \(\xi\). The additional first-order argument is what handles an
arbitrary specified finite action without assuming that its
representation preserves a chosen square-block decomposition.

\subsection{2. Endomorphisms, the hull image, and the matrix orbit map}\label{SF-E848:s2}

For this section put \(M=M_b\), with square size \(q\). Define

\[
 \begin{split}
 K&=\operatorname{Hom}_R(M,A^q),& k&=\dim_{\mathbf C}K,\\
 a_L&=\dim_{\mathbf C}\{t\in\operatorname{Mat}_q(A):
                         t b_{ij}=0\text{ for all }i<j\},\\
 \mathcal B&=\operatorname{im}\bigl(
       \operatorname{End}_R(M)\longrightarrow\operatorname{Mat}_q(R)
                                      \bigr),&
 c&=\operatorname{length}_R(\operatorname{Mat}_q(R)/\mathcal B).
 \end{split}
 \tag{MC.6}
\]

The hull-image map in this definition is the map induced on the
canonical double dual. Its kernel is exactly \(K\), by factorization
through the finite torsion. Applying \(\operatorname{Hom}_R(-,A^q)\)
to (MC.1), and using the truncated Koszul resolution, gives the exact
sequence

\[
 0\longrightarrow\operatorname{Mat}_q(A)^4
    \longrightarrow K\longrightarrow\operatorname{Mat}_q(A)
       \xrightarrow{\ t\mapsto (t b_{ij})\ }
                     \operatorname{Mat}_q(A)^6.
 \tag{MC.7}
\]

Therefore \(k=4q^2\ell+a_L\).

Every homomorphism of \(I^q\) extends uniquely to its hull \(R^q\).
A pair of endomorphisms, namely \(t\) on \(A^q\) and the restriction
of \(r\in\operatorname{Mat}_q(R)\) on \(I^q\), lifts to an
endomorphism of the extension precisely when their pushed-forward and
pulled-back Yoneda classes agree. In the Koszul coordinates this is
the actual matrix equation \(t b_{ij}=b_{ij}\bar r\) for every pair,
where the bar denotes reduction modulo \(I\). This criterion is valid
for arbitrary homomorphisms, not just automorphisms. In particular
\(\operatorname{Mat}_q(I)\subset\mathcal B\), by taking \(t=0\).
Thus \(c\) is finite, and

\[
 \overline{\mathcal B}
   =\{r\in\operatorname{Mat}_q(A):
         \exists t\in\operatorname{Mat}_q(A),\ t b_{ij}=b_{ij}r
          \text{ for every }i<j\},
 \qquad
 \dim\overline{\mathcal B}=q^2\ell-c.
 \tag{MC.8}
\]

Here \(\overline{\mathcal B}=\mathcal B/\operatorname{Mat}_q(I)\).
Consider the linear map

\[
 \psi:\operatorname{Mat}_q(A)\oplus\operatorname{Mat}_q(A)
      \longrightarrow\operatorname{Mat}_q(A)^6,
 \qquad \psi(t,r)=(t b_{ij}-b_{ij}r)_{i<j}.
 \tag{MC.9}
\]

Projection of its kernel onto the second factor is surjective onto
\(\overline{\mathcal B}\). Its kernel consists of the left annihilators
counted by \(a_L\). Consequently

\[
 \dim\operatorname{im}\psi=q^2\ell+c-a_L,
 \qquad k=4q^2\ell+a_L.
 \tag{MC.10}
\]

No splitting of \(\operatorname{End}_R(M)\to\mathcal B\) is asserted.
In contrast with the scalar case, the image need not be all of the
matrix algebra. The correction \(c\) will be retained until it cancels.

\subsection{3. A relative matrix Koszul family and its exact tangent kernel}\label{SF-E848:s3}

Put \(D=\mathbf C[\varepsilon]/(\varepsilon^2)\) and
\(S=D[[x_1,x_2,x_3,x_4]]\). Fix complex-linear lifts of \(A\) to
\(R\). For parameters \(\alpha\in A^4\) and
\(\beta\in\operatorname{Mat}_q(A)^6\), set

\[
 f_i^\varepsilon=f_i+\varepsilon\widetilde\alpha_i,
 \quad I_\varepsilon=(f_1^\varepsilon,\ldots,f_4^\varepsilon),
 \quad A_\varepsilon=S/I_\varepsilon,
 \quad b_{ij}^\varepsilon=\widetilde b_{ij}
                   +\varepsilon\widetilde\beta_{ij}\pmod{I_\varepsilon}.
 \tag{MC.11}
\]

These are deformations over a fixed commutative algebra \(S\);
coordinate automorphisms of \(R\) are not included in the equivalence.
The short exact sequence of Koszul complexes
\(0\to K(f;R)\to K(f^\varepsilon;S)\to K(f;R)\to0\), with first
map multiplication by \(\varepsilon\), proves that the middle complex
is acyclic in negative degrees. It also gives
\(0\to A\to A_\varepsilon\to A\to0\), in which multiplication by
\(\varepsilon\) is the quotient followed by the inclusion. Hence
\(\ker\varepsilon=\varepsilon A_\varepsilon\), so
\(A_\varepsilon\) is finite free over \(D\) of rank \(\ell\).
This last implication follows directly by lifting a complex basis of
\(\varepsilon A_\varepsilon\). The kernel \(I_\varepsilon\) of
\(S\to A_\varepsilon\) is \(D\)-flat as well.

The self-dual Koszul resolution gives
\(\operatorname{Ext}^j_S(A_\varepsilon,S)=0\) for \(j<4\).
The truncated resolution identifies
\(\operatorname{Ext}^1_S(I_\varepsilon^q,A_\varepsilon^q)\) with
\(\operatorname{Mat}_q(A_\varepsilon)^6\). Thus (MC.11) defines an
actual \(D\)-flat extension \(M_\varepsilon\) reducing to \(M\).
One explicit resolution uses the two Koszul complexes with triangular
differential \(\left[\begin{smallmatrix}d&B_\varepsilon\\0&-d\end{smallmatrix}\right]\),
where \(B_\varepsilon\) is contraction by the matrix-valued two-form
\(b^\varepsilon\). Even contraction commutes with the scalar Koszul
differential. No products of two \(B_\varepsilon\) blocks occur in
the square, so noncommutation of the six coefficient matrices causes
no problem. The short exact sequence of these complexes verifies the
extension and its exact resolution directly.

The relative hull is also explicit over the nonreduced ring \(S\):

\[
 \operatorname{Hom}_S(M_\varepsilon,S)
       =\operatorname{Hom}_S(I_\varepsilon^q,S)=S^q,
 \qquad M_\varepsilon^{**}=S^q,
 \qquad
 \operatorname{im}(M_\varepsilon\to M_\varepsilon^{**})
       =I_\varepsilon^q.
 \tag{MC.12}
\]

These equalities follow from the just proved grade-four vanishings.
Their generators are the actual coordinate quotient-to-hull maps,
so the identifications commute with reduction to \(R\). Neither
regularity of \(S\) nor a general double-dual base-change assertion
has been used.

The family gives a complex-linear tangent map
\(L:A^4\oplus\operatorname{Mat}_q(A)^6\to
\operatorname{Ext}^1_R(M,M)\). Differentiation of its actual free
differential gives the linear cocycle. A zero tangent class is
equivalent to an \(S\)-linear trivialization of the module deformation
with identity special fiber: a splitting of the self-extension
\(0\to M\to M_\varepsilon\to M\to0\) supplies such a trivialization
by \(m_0+\varepsilon m_1\mapsto s(m_0)+\varepsilon s(m_1)\), and the
converse follows by restriction.

By (MC.12), a trivialization induces a matrix in \(GL_q(S)\) sending
\(I_\varepsilon^q\) to \((I\otimes D)^q\). Every invertible matrix
over \(S\) preserves \(J^q\) for every ideal \(J\subset S\): both
it and its inverse have entries in \(S\). Therefore these two ideal
submodules are equal, and hence their coordinate ideals are equal.
It follows that every \(\alpha_i\) is zero in \(A\).

With \(\alpha=0\), the induced torsion and hull automorphisms are
\(1+\varepsilon t\) and \(1+\varepsilon r\), respectively. The
identity of their extension classes is
\((1+\varepsilon t)(b+\varepsilon\beta)
 =b(1+\varepsilon\bar r)\). Thus
\(\beta=b\bar r-tb\). Conversely any tuple of this form is killed
by these actual changes of torsion and hull trivialization, using
Yoneda equivalence of extensions. Every matrix over \(A\) lifts to
one over \(R\). We have proved the exact kernel and the inequality

\[
 \ker L=\{0\}\oplus\operatorname{im}\psi,
 \qquad
 e_1\ge(4+6q^2)\ell-(q^2\ell+c-a_L)
          =(4+5q^2)\ell-c+a_L.
 \tag{MC.13}
\]

Only this explicit family has been counted. There is no assertion that
every deformation of \(M\) preserves its common-CI or matrix form.

\subsection{4. The finite-supported cone and the exact middle-Ext identity}\label{SF-E848:s4}

Let \(C=\operatorname{cone}(M\xrightarrow{\phi}R^q)\), with
cohomology \(A^q\) in degrees \(-1\) and \(0\). Both \(M\) and
\(C\) are perfect. The complex \(F=R\operatorname{Hom}_R(M,C)\)
has finite-length cohomology. A finite free resolution of \(M\)
gives its Euler characteristic without subtracting infinite dimensions:

\[
 \chi_{\mathrm{length}}(F)
   =\left(\sum_i(-1)^i\operatorname{rank}L_i\right)
          \chi_{\mathrm{length}}(C)=0.
 \tag{MC.14}
\]

Write \(h_j=\dim\operatorname{Ext}^j_R(M,C)\). The grade-four
vanishings for the two finite defects give
\(\operatorname{Ext}^1_R(M,R^q)=\operatorname{Ext}^2_R(M,R^q)=0\).
Applying \(\operatorname{Hom}_R(M,-)\) to the hull triangle gives
\(h_{-1}=k\), \(h_0=e_1+c\), and \(h_1=e_2\). In the middle
identity the exact sequence is
\(0\to\operatorname{Mat}_q(R)/\mathcal B\to
\operatorname{Hom}_R(M,C)\to\operatorname{Ext}^1_R(M,M)\to0\).

Finite-support local duality and perfect tensor-Hom adjunction give

\[
 \operatorname{Ext}^j_R(M,C)^\vee
          \simeq\operatorname{Ext}^{4-j}_R(C,M).
 \tag{MC.15}
\]

The dual is the complex-vector-space dual, using a coordinate volume
for the residue identification. This is applied only to \(F\), whose
cohomology has finite length. In detail
\(R\operatorname{Hom}(F,R[4])
 \simeq R\operatorname{Hom}(C,M)[4]\); the normalized dualizing
complex of \(R\) is \(R[4]\). This is the local-duality input of
Stacks Section 47.18, tag 0A81, Theorem 47.18.3 and Lemma 47.18.4,\cite{SFStacksE847LocalDuality}
with the residue identification in Lipman Lemma 7.3/Theorem 7.4,
pp.~64--69,\cite{SFLipman1984} as explicitly invoked and checked in E847. It does not
discard local cohomology of \(\operatorname{End}_R(M)\).

For completeness compute the three groups on the right needed here.
The kernel of \(\operatorname{Hom}_R(R^q,M)\to
\operatorname{End}_R(M)\) consists of the maps into \(A^q\).
Indeed a map in that kernel vanishes on \(I^q\); its columns are
annihilated by \(I\). The annihilator of \(I\) in \(M\) is exactly
\(A^q\), since the quotient \(I^q\) has no such nonzero element
over the domain \(R\). Consequently
\(\operatorname{Hom}_R(C,M)\) has dimension \(q^2\ell\).

The image of \(\operatorname{Hom}_R(R^q,M)\) in
\(\operatorname{End}_R(M)\) projects isomorphically onto
\(\operatorname{Mat}_q(I)\). Surjectivity follows by lifting the
columns of a matrix in \(I^q\) through \(M\to I^q\). If such
an image has zero hull matrix, all lifted columns lie in \(A^q\);
their composite with \(\phi\) is zero because \(I A=0\).
Thus this image has zero intersection with \(K\). The same hull
triangle now gives

\[
 \begin{split}
 \dim\operatorname{Hom}_R(C,M)&=q^2\ell,\\
 0\longrightarrow K\longrightarrow\operatorname{Ext}^1_R(C,M)
      &\longrightarrow\mathcal B/\operatorname{Mat}_q(I)
                      \longrightarrow0,\\
 \operatorname{Ext}^2_R(C,M)&\simeq\operatorname{Ext}^1_R(M,M).
 \end{split}
 \tag{MC.16}
\]

The last identity uses \(\operatorname{Ext}^{>0}_R(R^q,M)=0\).
Equations (MC.15)--(MC.16) and the earlier terms therefore give all
six finite cohomology dimensions:

\[
 (h_{-1},h_0,h_1,h_2,h_3,h_4)
   =(k,\ e_1+c,\ e_2,\ e_1,\ k+q^2\ell-c,\ q^2\ell).
 \tag{MC.17}
\]

Substitution in (MC.14) proves the exact identity, and then the bound:

\[
 \begin{split}
 e_2&=2e_1-2k+2c,\\
 e_2&\ge2\bigl((4+5q^2)\ell-c+a_L\bigr)
                -2(4q^2\ell+a_L)+2c
       =(2q^2+8)\ell.
 \end{split}
 \tag{MC.18}
\]

This proves SF-E848a. Both the hull-image correction and the left
annihilator cancel. No commutation, invertibility, or nonvanishing
condition on the tuple \(b\) was used.

\subsection{5. Actual stabilizer summands may initially be rectangular}\label{SF-E848:s5}

Assume (MC.3)--(MC.4), and consider the stabilizer \(G_p\) of the
point under the specified translation/tensor action. A translation
fixing \(p\) is the identity translation. Hence this stabilizer acts
\(R\)-linearly on the completed module. Its Picard factors are
torsion line bundles. In flat local frames their coherent multiplication
isomorphisms have constant factors, giving one projective scalar
cocycle \(\alpha\in Z^2(G_p,\mathbf C^*)\). This is precisely the
point-stabilizer interface of E846--E847 and supplied S21, not an
assertion about arbitrary finite automorphisms of the formal base.

The twisted finite group algebra \(\mathbf C_\alpha[G_p]\) is
semisimple. Averaging conjugates of a projection proves the usual
Maschke argument; projective scalar factors cancel under conjugation.
Let \(S_\lambda\) be its simple modules. There is an actual
\(R\)-linear decomposition

\[
 \widehat E_p=\bigoplus_\lambda S_\lambda\otimes N_\lambda,
 \qquad
 N_\lambda=\operatorname{Hom}_{\mathbf C_\alpha[G_p]}
                         (S_\lambda,\widehat E_p).
 \tag{MC.19}
\]

The action is on \(S_\lambda\) alone. We do not assume a conjugacy
theorem for an action on a preselected decomposition of \(M_b\).

Apply the exact multiplicity functor in (MC.19) to the canonical
torsion, hull, image, and quotient sequences. The original torsion
and quotient at \(p\) are both \(A^m\). Their multiplicity modules
are direct summands of finite direct sums of \(A^m\): the appropriate
twisted-group averaging idempotent is \(A\)-linear. They are therefore
finite projective, hence free, over the local algebra \(A\). Similarly
the hull multiplicity is free over \(R\). Write their ranks as
\(t_\lambda,q_\lambda,h_\lambda\), respectively.

The surjection \(R^{h_\lambda}\to A^{q_\lambda}\) has a unit
\(q_\lambda\) by \(q_\lambda\) minor after reduction modulo
\(\mathfrak m\). Elementary row operations over \(A\), and column
operations lifted to \(R\), put it in the form of the residue map
on the first \(q_\lambda\) coordinates and zero on the remaining
coordinates. Its kernel is consequently
\(I^{q_\lambda}\oplus R^{h_\lambda-q_\lambda}\).
The extension over the free summand splits, because
\(\operatorname{Ext}^1_R(R,A^{t_\lambda})=0\). Thus

\[
 N_\lambda\simeq M_{b_\lambda}\oplus R^{h_\lambda-q_\lambda},
 \quad
 0\longrightarrow A^{t_\lambda}\longrightarrow M_{b_\lambda}
      \longrightarrow I^{q_\lambda}\longrightarrow0,
 \quad
 b_\lambda\in\operatorname{Mat}_{t_\lambda\times q_\lambda}(A)^6.
 \tag{MC.20}
\]

This is an actual splitting on the multiplicity module. Any such
choice is still equivariant after tensoring with \(S_\lambda\),
because the group acts on that representation factor alone. At this
point the two ranks in (MC.20) have not been assumed equal. Establishing
their equality is essential; an inequality for square matrices cannot
simply be applied to a rectangular isotypic component.

\subsection{6. First-order evaluation forces equality in each summand}\label{SF-E848:s6}

Assume the vanishing hypothesis of SF-E848b. The characteristic map
is natural for restriction, completion, and direct sums. The higher
HKR components \(H^1(T_X)\) and \(H^2(\mathcal O_X)\) restrict to
zero on a formal affine neighborhood. In coordinates supplied by the
abelian formal logarithm, the remaining bivector is a nonzero constant
bivector \(\pi\). Thus \(\rho_\pi(\widehat E_p)=0\), and diagonal
restriction through the actual decomposition (MC.19) gives
\(\rho_\pi(N_\lambda)=0\) for every \(\lambda\).

Here is the exact local implication being used, including the needed
finite-support trace. The first-order algebra
\(R\oplus\varepsilon R\) has product
\(a*b=ab+\varepsilon\pi(da,db)/2\). It is associative modulo
\(\varepsilon^2\) by the Leibniz rule. The obstruction to deforming
an \(R\)-module over it is its Hochschild characteristic class in
\(\operatorname{Ext}^2\). A finite free resolution can be lifted
coefficientwise; its squared differential is this cocycle, and a
degree-one correction kills it when the class is zero. The reduced
complex resolves the module, so the square-zero cohomology sequence
gives a flat first-order module deformation. This is the continuous
finite-resolution argument of E819 Section 3; no discontinuous action
on formal power series or all-order deformation is used.

For a module \(N\) with finite kernel \(T_N\) and cokernel \(Q_N\)
in a free hull map \(N\to R^h\), grade four gives
\(\operatorname{Ext}^1_R(N,R^h)=0\). Hence the hull map lifts to the
first-order deformed module and the free deformed hull. This follows
by applying the Hom complex to lifted finite free resolutions: the
obstruction to a lifted degree-zero cocycle lies in that vanishing
Ext group. The defect cone is a two-term complex of flat dual-number
modules with finite cohomology. It is perfect over the dual numbers,
by finite Tor amplitude and bounded finite cohomology. The exact
criterion is Stacks Lemma 15.76.2, tag 0658,\cite{SF-StacksPerfect} as invoked in E819 Section 2.2.

On a bounded finite free representative of this cone, the supertrace
of a commutator vanishes, as does that of a null-homotopic map. The
coordinate operators satisfy
\([x_i,x_j]=\varepsilon\pi_{ij}\operatorname{Id}\). Transferring
them to that perfect representative therefore gives

\[
 0=\varepsilon\pi_{ij}
             \bigl(\operatorname{length}Q_N
                         -\operatorname{length}T_N\bigr).
 \tag{MC.21}
\]

Some \(\pi_{ij}\ne0\), and \(\varepsilon\ne0\), so the two
lengths agree. The trace here is on the finite perfect defect cone,
not on the infinite-dimensional free hull.

Apply this argument to each actual \(N_\lambda\). Its canonical
finite defects have lengths \(t_\lambda\ell\) and
\(q_\lambda\ell\). It follows that
\(t_\lambda=q_\lambda\). If \(q_\lambda=0\), this also proves
\(t_\lambda=0\), and that component is free. Since the original
quotient \(A^m\) is nonzero, some \(q_\lambda>0\). For that
component (MC.20) is a square matrix extension plus a free module.
No equivariant noncommutative deformation is required: ordinary
evaluation vanishes on the actual summand and the local argument is
ordinary. Nor was a balanced character on the summand asserted.

\subsection{7. The global invariant local-Ext space and its trace rank}\label{SF-E848:s7}

We include the actual injection and its trace factorization to avoid
replacing a global connecting map by an unsupported localization
claim. Put \(\mathcal C=\operatorname{cone}(E\to H)\). All of its
cohomology is supported on a finite set \(Z\). Since \(E\) is perfect
on the smooth variety \(X\), the cohomology sheaves of
\(R\mathcal Hom(E,\mathcal C)\) are finite-supported. Their higher
sheaf cohomology vanishes. Restriction and completion consequently give

\[
 \operatorname{Ext}^j_X(E,\mathcal C)
   \simeq\bigoplus_{p'\in Z}
       \operatorname{Ext}^j_{\widehat{\mathcal O}_{X,p'}}(M_{p'},C_{p'}).
 \tag{MC.22}
\]

At every point, the grade-four finite-defect argument gives
\(\operatorname{Ext}^1_R(M_{p'},R^{h})
 =\operatorname{Ext}^2_R(M_{p'},R^{h})=0\). The local hull triangle
identifies \(\operatorname{Ext}^1_R(M_{p'},C_{p'})\) with
\(\operatorname{Ext}^2_R(M_{p'},M_{p'})\). The global preceding
arrow \(\operatorname{Ext}^1_X(E,H)\to
\operatorname{Ext}^1_X(E,\mathcal C)\) is zero: in the natural
restriction square its bottom source is zero, whereas its right
vertical arrow is the isomorphism (MC.22). The global connecting
homomorphism is therefore a natural injection

\[
 j:\bigoplus_{p'\in Z}
          \operatorname{Ext}^2_R(M_{p'},M_{p'})
       \hookrightarrow\operatorname{Ext}^2_X(E,E).
 \tag{MC.23}
\]

It is equivariant for the actual finite action because the hull triangle
is canonical. It is also a section of restriction to these local Ext
groups, with the stated local identifications.

For a local class \(u\), let \(a\in\operatorname{Ext}^1_X(E,\mathcal C)\)
be its preimage under these identifications. For a degree-two Hochschild
class \(\gamma\), associativity of composition and the connecting
map express the Serre pairing of \(j(u)\) with \(\rho_E(\gamma)\)
as the trace of the image of \(a\rho_E(\gamma)\) in
\(\operatorname{Ext}^3_X(E,\mathcal C)\), up to the fixed harmless
cyclic sign. By (MC.22), that product depends only on the actual
completed restrictions of the evaluation. The HKR components
\(H^1(T_X)\) and \(H^2(\mathcal O_X)\) have zero restriction,
so the pairing factors through

\[
 HH^2(X)\longrightarrow H^0(X,\Lambda^2T_X),
 \qquad \dim_{\mathbf C}H^0(X,\Lambda^2T_X)=6.
 \tag{MC.24}
\]

The factorization through the finite-support Ext group, rather than
only a statement about locally vanishing classes, excludes hidden
hypercohomology-boundary contributions. The supplied CY4 Serre
trace/evaluation adjunction identifies the annihilator of all these
evaluations with the kernel of the full semiregularity map. Its exact
scope is DEG\_L S20, first proof paragraph, and S21.6 through the
adjunction immediately preceding S21.8. The ordinary Serre pairing is
Stacks 48.27, Lemma 48.27.1(5),(7), Remarks 48.27.2--3 and Lemma 48.27.4.\cite{SFStacksE846Duality}
This uses that adjunction, not merely the composition formula
\(\sigma\rho=\operatorname{contraction}\operatorname{ch}\).
Thus on any \(d\)-dimensional subspace of the source of (MC.23),
the semiregularity kernel has dimension at least \(d-6\).

Choose the component \(\lambda\) with \(q_\lambda>0\) found in
Section 6. For any
\(u\in\operatorname{Ext}^2_R(M_{b_\lambda},M_{b_\lambda})\),
place \(\operatorname{Id}_{S_\lambda}\otimes u\) in that actual
summand of the local module and zero in all its other blocks, including
the free block. This is a \(G_p\)-invariant class. It is nonzero for
nonzero \(u\), so it defines an invariant space of dimension
\(e_2(M_{b_\lambda})\). Transport over the orbit of \(p\) and sum
the components. Stabilizer invariance makes this independent of coset
representatives, and its component at \(p\) proves injectivity.
The canonical map (MC.23) takes it into invariant global Ext.

The six-dimensional trace bound (MC.24) still applies after this orbit
sum, since the higher HKR components vanish at every point. In the
specified effective translation/tensor framework, supplied S21 proves
that the other sectors vanish and identifies this ordinary map on
invariant Ext with the relevant identity-sector semiregularity map.
We have obtained

\[
 \dim\ker\bigl(\sigma_E^2|_{\operatorname{Ext}^2(E,E)^G}\bigr)
   \ge (2q_\lambda^2+8)\ell-6
   \ge10\ell-6>0.
 \tag{MC.25}
\]

This proves SF-E848b. When the action is the identity, one may apply
the same argument directly to the given square block (MC.4); Sections5
and6 are unnecessary, and no Hochschild-evaluation vanishing hypothesis
is needed in that special case.

\subsection{8. Why the balanced character supplies the vanishing hypothesis}\label{SF-E848:s8}

We repeat the precise accepted E819 interface. For each RM plane
choose a nonzero invariant holomorphic bivector \(\pi_\pm\) and
define \(\eta_\pm\in H^2(\mathcal O_X)\) by
\(\pi_\pm\lrcorner(\theta_\pm^2/2)=\eta_\pm\). Contraction
by \(\pi_\pm\) on its factor in (MC.5) gives
\(-\nu\eta_\pm\), while wedge by \(\nu\eta_\pm\) gives the
opposite term; wedge with the top term of that same plane is zero.
Hence \(\pi_\pm+\nu\eta_\pm\) annihilates the full balanced
product under the Hochschild cap action. The divisor shear for \(B\)
preserves the bivector component, as follows by expanding
\(e^{-B}\iota_v e^B=\iota_v+(\iota_vB)\wedge\). We obtain
\(\xi\in\ker c^2_{\operatorname{ch}(E)}\) with nonzero bivector
component, for the entire character including the rational B-field.

Published Buchweitz--Flenner Theorem 6.4.2 and Proposition 6.2.1\cite{SF-BF2008} identify
semiregularity after evaluation with contraction by the full Chern
character. Supplied S21 gives its compatibility with the coherent action
and ensures that the evaluation takes values in invariant Ext. If
\(\sigma^2_{E,G}\) were injective, then
\(\sigma^2_{E,G}\rho^2_{E,G}(\xi)=0\) would imply
\(\rho^2_{E,G}(\xi)=0\). Its ordinary Ext image is zero, which is
exactly the hypothesis of SF-E848b. The resulting contradiction to
(MC.25) proves SF-E848c. Only a first-order deformation implication
was used; no second-order or formal noncommutative unobstructedness
claim appears.

\subsection{9. Actual realizability, source status, and limitations}\label{SF-E848:s9}

The theorem concerns actual local modules, and every displayed matrix
tuple can occur in an actual coherent sheaf. Since \(I\) has finite
colength, the algebra \(A\) is a quotient of a finite jet of the
algebraic local ring of \(X\). It defines a finite subscheme \(Z\).
The finite-support spectral sequence identifies
\(\operatorname{Ext}^1_X(I_Z^q,\mathcal O_Z^q)\) with its completed
local value \(\operatorname{Mat}_q(A)^6\). Thus each \(b\) defines
an actual extension with hull \(\mathcal O_X^q\), torsion and
quotient \(\mathcal O_Z^q\), and Chern character \(q\).

Balanced coherent examples are obtained as follows, without asserting
semiregularity. Take the actual rank 32 balanced hull \(H_0\) of
E815/E820 and its direct sum \(H_0^{\oplus q}\). A choice of a local
frame gives a surjection from this hull onto \(\mathcal O_Z^q\).
Its kernel has completed module \(I^q\oplus R^{31q}\). Every matrix
class \(b\) lifts to an actual extension of that kernel by
\(\mathcal O_Z^q\), because its relevant Ext has finite support and
the free local summand has no Ext-one to that finite module. The
resulting sheaf has the local type (MC.4), and its torsion and hull
quotient cancel in its full Chern character. The identity group is
an actual action. No additional finite action or indecomposability is
claimed for this construction.

The universal matrix calculation, the relative canonical hull, the
exact tangent kernel, and the actual isotypic reduction are written in
full. The matrix argument has shared root/sixfolds provenance. The
first-order finite-cone lemma and its cap-kernel application are
accepted E819 inputs, whose Sections 2--4 were freshly reread in full
for this proof. The local/global duality and support interfaces are
those just checked in E847; no fresh numerical certificate or repeated
whole reading of every upstream source is implied. The invoked
published Buchweitz--Flenner locators are Proposition 6.2.1,
pp.~29--30; Proposition 6.3.1 and Corollary 6.3.2, pp.~30--31;
Theorem 6.4.2 and proof through Proposition 6.4.4, pp.~31--33 of the
author copy of Advances in Mathematics 217 (2008),243--281,
DOI 10.1016/j.aim.2007.06.013. Supplied S20/S21 retains its supplied
proved-input status throughout.

The common complete-intersection algebra at the chosen point is
essential to this statement. Distinct finite torsion and quotient
algebras, arbitrary finite modules over one ideal, nonfree hulls,
positive-dimensional defects, general derived objects, and actions
with nontrivial formal base automorphisms are not covered. The global
unconditional finite-action statement of scalar E847 is not silently
extended: the matrix theorem uses the explicit evaluation hypothesis,
or the balanced character to force it. A concrete comparison with all
earlier local criteria is still required before any additional strict
main-gate movement is asserted.
